\documentclass[12pt]{article}

\usepackage[utf8]{inputenc}
\usepackage[spanish]{babel}
\usepackage[shortlabels]{enumitem}

\usepackage{mathtools}
\usepackage{amssymb}
\usepackage{amsthm}

\usepackage{geometry}
\geometry{margin=2.5cm}

\usepackage{hyperref}

\title{Ejercicios Espacios Vectoriales}
\author{}

\begin{document}
\maketitle
\section*{Ejercicio 1}

En el espacio vectorial \(P_3(x)\) de los polinomios de grado menor o igual que \(3\), se consideran los subespacios
\[
S_1=\left\{p(x)\in P_3(x)\;:\; p(0)=0 \text{ y las tangentes a } p(x) \text{ en los puntos de abscisas } 1 \text{ y } -1 \text{ son paralelas}\right\},
\]
\[
S_2=\left\{p(x)\in P_3(x)\;:\; p(2)=0\right\}.
\]

\begin{enumerate}
    \item Calcular la dimensión y obtener una base de cada uno de esos dos subespacios.
    \item Calcular el subespacio \(S_1\cap S_2\), una base del mismo y razonar si \(S_1+S_2\) es o no una suma directa.
    \item Si \(p(x)\) es un vector de \(S_1\cap S_2\), calcular cuáles son sus coordenadas en la base
    \[
    B=\{1+x,\;x+x^2,\;x^2+x^3,\;x^3\}.
    \]
    \item Demostrar, utilizando la matriz de cambio de base apropiada, que las coordenadas en la base canónica del vector obtenido en el apartado \textbf{c)} son efectivamente las mismas.
\end{enumerate}

\section*{Ejercicio 2}

Calcular la dimensión y una base del subespacio
\[
U=\left\{X\in M_{2\times 2}(\mathbb{R}) \;:\; \operatorname{Traza}(MX)=0\right\},
\qquad
\text{con }
M=
\begin{pmatrix}
1 & -1\\
-1 & 1
\end{pmatrix}.
\]


\section*{Ejercicio 3}

Dados los subespacios de \(\mathbb{R}^4\):
\[
U=\left\langle (1,0,0,1),\,(1,1,1,1),\,(0,2,2,0)\right\rangle
\]
y
\[
W=\left\{(x,y,z,t)\in\mathbb{R}^4 \;:\; 3x+y-z-3t=0,\; y-t=0\right\}.
\]

Determinar razonadamente si \(U+W\) es una suma directa.

\section*{Ejercicio 4}

Si \(U\) y \(V\) son subespacios vectoriales de \(\mathbb{R}^{100}\) tales que
\[
\dim U=60
\qquad\text{y}\qquad
\dim V=50,
\]
razonar si las siguientes afirmaciones son verdaderas o falsas:

\begin{enumerate}
    \item Un sistema generador de \(U\cap V\) puede estar formado por \(9\) vectores.
    \item El espacio vectorial \(V\) puede generarse con \(55\) vectores. En caso de respuesta negativa, razonarlo y en caso de respuesta afirmativa, indicar cómo se construiría un sistema generador.
\end{enumerate}

\section*{Ejercicio 5}

Si \(P_2(x)\) es el espacio vectorial de los polinomios de grado menor o igual que \(2\) y \(p(x)\) es un polinomio de grado exactamente \(2\),

\begin{enumerate}
    \item[a)] Demostrar que tanto
    \[
    B=\{p(x),\,p'(x),\,p''(x)\}
    \]
    como
    \[
    B'=\{p(x),\,p(x)+p'(x),\,p'(x)+p''(x)\}
    \]
    son bases de \(P_2(x)\).

    \item[b)] Analizar si
    \[
    S=\{x(x-a)\,:\,a\in\mathbb{R}\}
    \]
    es un subespacio vectorial de \(P_2(x)\).

    \item[c)] Si \(p(x)\in S\) y tiene raíz \(-1\), obtener las coordenadas de
    \[
    q(x)=x^2+x+2
    \]
    en la base \(B'\).

    \item[d)] Demostrar, utilizando la matriz de cambio de base apropiada, que el vector obtenido en c) es precisamente \(q(x)\).
\end{enumerate}

\section*{Ejercicio 6}
Sea \(A=\{v_1,v_2,v_3,v_4,v_5\}\) un subconjunto de un espacio vectorial \(E\) con \(\dim E=4\). Razonar si las siguientes afirmaciones son verdaderas o falsas:
    
\begin{enumerate}
    \item Si \(A\) genera \(E\), entonces de \(A\) puede extraerse una base de \(E\).
    \item Si todo subconjunto de \(4\) vectores de \(A\) es linealmente dependiente, entonces \(A\) no genera \(E\).
    \item Si \(A\) contiene una base de \(E\), entonces \(A\) genera \(E\).
    \item Si \(A\) es linealmente dependiente, entonces \(A\) no puede contener una base de \(E\).
\end{enumerate}

\end{document}